\documentclass{article}
\usepackage{amsmath, amssymb, amsthm}

\title{IMC 2025 \\ Second Day, July 31, 2025}
\date{}

\begin{document}
\maketitle

\section*{Problem 6}
Let $f : (0, \infty) \to \mathbb{R}$ be a continuously differentiable function, and let $b > a > 0$ be real numbers such that $f(a) = f(b) = k$. Prove that there exists a point $\xi \in (a,b)$ such that
\[
f(\xi) - \xi f'(\xi) = k.
\]
(10 points)

\section*{Problem 7}
Let $\mathbb{Z}_{>0}$ be the set of positive integers. Find all nonempty subsets $M \subseteq \mathbb{Z}_{>0}$ satisfying both of the following properties:
\begin{itemize}
\item[(a)] if $x \in M$, then $2x \in M$,
\item[(b)] if $x, y \in M$ and $x + y$ is even, then $\dfrac{x+y}{2} \in M$.
\end{itemize}
(10 points)

\section*{Problem 8}
For an $n \times n$ real matrix $A \in M_n(\mathbb{R})$, denote by $A^{\mathsf{R}}$ its counter-clockwise $90^\circ$ rotation. For example,
\[
\begin{bmatrix} 1 & 2 & 3 \\ 4 & 5 & 6 \\ 7 & 8 & 9 \end{bmatrix}^{\mathsf{R}} = \begin{bmatrix} 3 & 6 & 9 \\ 2 & 5 & 8 \\ 1 & 4 & 7 \end{bmatrix}.
\]
Prove that if $A = A^{\mathsf{R}}$ then for any eigenvalue $\lambda$ of $A$, we have $\operatorname{Re}\lambda = 0$ or $\operatorname{Im}\lambda = 0$.

(10 points)

\section*{Problem 9}
Let $n$ be a positive integer. Consider the following random process which produces a sequence of $n$ distinct positive integers $X_1, X_2, \ldots, X_n$.

First, $X_1$ is chosen randomly with $\mathbb{P}(X_1 = i) = 2^{-i}$ for every positive integer $i$. For $1 \leq j \leq n-1$, having chosen $X_1, \ldots, X_j$, arrange the remaining positive integers in increasing order as $n_1 < n_2 < \cdots$, and choose $X_{j+1}$ randomly with $\mathbb{P}(X_{j+1} = n_i) = 2^{-i}$ for every positive integer $i$.

Let $Y_n = \max\{X_1, \ldots, X_n\}$. Show that
\[
\mathbb{E}[Y_n] = \sum_{i=1}^n \frac{2^i}{2^i - 1}
\]
where $\mathbb{E}[Y_n]$ is the expected value of $Y_n$.

(10 points)

\section*{Problem 10}
For any positive integer $N$, let $S_N$ be the number of pairs of integers $1 \leq a, b \leq N$ such that the number $(a^2 + a)(b^2 + b)$ is a perfect square. Prove that the limit
\[
\lim_{N \to \infty} \frac{S_N}{N}
\]
exists and find its value.

(10 points)

\end{document}
